\chapter{Implementação}
\label{cap:implementacao}

O sistema foi implementado utilizando a linguagem C++17 e gerenciado por um \texttt{Makefile} unificado na raiz do projeto. A aplicação \textbf{IBE Mail} funciona como um cliente de e-mail integrado a um gerador de chaves privadas (PKG), executando localmente sem necessidade de servidor externo.

\section{Bibliotecas Utilizadas}
\label{sec:bibliotecas}

O desenvolvimento apoiou-se no seguinte conjunto de bibliotecas e ferramentas:

\begin{itemize}
    \item \textbf{PBC (Pairing-Based Cryptography):} Biblioteca central do projeto, que provê as operações de emparelhamento bilinear sobre curvas elípticas supersingulares necessárias para o algoritmo de Boneh-Franklin.
    \item \textbf{GMP (GNU Multiple Precision Arithmetic):} Utilizada como dependência fundamental da PBC para o tratamento de números inteiros de precisão arbitrária.
    \item \textbf{OpenSSL (\texttt{libcrypto}):} Responsável por fornecer a função criptográfica de dispersão SHA-256, a qual foi adaptada com separação de domínio para modelar os quatro oráculos aleatórios ($H_1$ a $H_4$).
    \item \textbf{Dear ImGui:} Biblioteca C++ de interface gráfica imediata (\textit{Immediate Mode GUI}), inserida diretamente no código-fonte do projeto.
    \item \textbf{SDL2 e OpenGL 3:} Bibliotecas de baixo nível responsáveis pela criação da janela no sistema operacional, captura de eventos e aceleração por \textit{hardware} para a renderização gráfica do ImGui.
\end{itemize}

\section{Estrutura do Projeto}
\label{sec:estrutura}

A arquitetura do código está segmentada em módulos funcionais distintos dentro do diretório \texttt{implementation/}, garantindo a separação entre a matemática, a lógica de negócio e a interface:

\begin{itemize}
    \item \textbf{\texttt{src/common/}}: Contém as primitivas criptográficas puras. O arquivo \texttt{BonehFranklin} implementa os algoritmos \textit{Setup}, \textit{Extract}, \textit{Encrypt} e \textit{Decrypt}. O arquivo \texttt{Hash} envelopa as chamadas do OpenSSL, enquanto \texttt{Utils} manipula UUIDs, formatação Base64 e operações XOR.
    
    \item \textbf{\texttt{src/core/}}: Hospeda a lógica de negócio da aplicação. O módulo \texttt{PKG} gerencia o ciclo de vida dos usuários (registro, aprovação, extração e revogação de chaves), interagindo com o arquivo \texttt{registry.dat}. O módulo \texttt{MailSystem} abstrai a estrutura do disco, sendo encarregado de escrever e ler os arquivos \texttt{.ibemail} cifrados, roteando mensagens entre as pastas de recebidos (\texttt{inbox/}) e enviados (\texttt{sent/}).
    
    \item \textbf{\texttt{src/gui/}}: Envolve a classe \texttt{App}, que consome os serviços do núcleo do sistema para alimentar os painéis visuais do programa (Login, Caixa de Entrada, Envio de E-mails e Painel Administrativo). 
    
    \item \textbf{\texttt{data/}}: Pasta criada em tempo de execução que atua como banco de armazenamento físico. Nela residem os parâmetros públicos (\texttt{params.pub}), o segredo mestre isolado do PKG (\texttt{master.sec}), a base de identidades, as chaves privadas extraídas (\texttt{keys/}) e as estruturas de diretório de e-mail.
\end{itemize}
